% -----------------------------------------------------------------------------
\subsubsection{Présence (F5)}
\label{sec:presence}

Tado règle le chauffage selon que la maison est habitée ou non. L'utilisateur peut laisser
Tado en décider d'après ses appareils mobiles (mode \emph{Auto}), ou l'imposer lui-même.

\begin{exigence}{TPB-PRE-001}{Présence actuelle}[
    justification = {Savoir ce que Tado croit, avant de le corriger.},
    priorite      = {Essentielle},
    verification  = {T},
    phase         = {Présence},
  ]
  L'application doit afficher, dans l'écran de présence, l'état de présence actuel de la
  maison (\emph{Présent} ou \emph{Absent}) et le mode qui le détermine.
\end{exigence}

\begin{exigence}{TPB-PRE-002}{Trois modes}[
    justification = {Forcer l'un ou l'autre état, ou laisser faire Tado.},
    priorite      = {Essentielle},
    verification  = {T},
    phase         = {Présence},
  ]
  L'application doit permettre à l'utilisateur de choisir entre trois modes de présence~:
  \emph{Présent}, \emph{Absent} et \emph{Auto}.
\end{exigence}

\begin{exigence}{TPB-PRE-003}{Mode Auto expliqué}[
    justification = {Comprendre pourquoi la maison est considérée comme présente ou
                     absente.},
    priorite      = {Importante},
    verification  = {T},
    phase         = {Présence},
    note          = {Tado décide d'après la position des téléphones qu'il suit (géolocalisation),
                     non d'après leur connexion au wifi~: c'est le choix retenu.},
  ]
  Lorsque le mode \emph{Auto} est choisi, l'application doit indiquer sur quoi repose
  l'état de présence, en donnant le nombre d'appareils mobiles présents à la maison.
\end{exigence}

\begin{exigence}{TPB-PRE-004}{Mode actuel repéré}[
    justification = {Voir le mode choisi sans le relire dans une phrase.},
    priorite      = {Importante},
    verification  = {D},
    phase         = {Présence},
  ]
  L'application doit marquer, parmi les trois modes, celui qui est en vigueur.
\end{exigence}

\begin{exigence}{TPB-PRE-005}{Changement confirmé}[
    justification = {Un chauffage laissé en absence par erreur coûte une soirée froide.},
    priorite      = {Essentielle},
    verification  = {T},
    phase         = {Présence},
  ]
  Après un changement de mode de présence, l'application doit indiquer s'il a été pris en
  compte et afficher le mode et l'état de présence qui en résultent.
\end{exigence}
